%% Created by Maple 2022.1, Windows 10
%% Source Worksheet: lab_sols07
%% Generated: Wed May 08 10:35:54 BST 2024
\documentclass{article}
\usepackage{amssymb}
\usepackage{graphicx}
\usepackage{hyperref}
\usepackage{listings}
\usepackage{mathtools}
\usepackage{maple}
\usepackage[utf8]{inputenc}
\usepackage[svgnames]{xcolor}
\usepackage{amsmath}
\usepackage{breqn}
\usepackage{textcomp}
\begin{document}
\lstset{basicstyle=\ttfamily,breaklines=true,columns=flexible}
\pagestyle{empty}
\DefineParaStyle{Maple Bullet Item}
\DefineParaStyle{Maple Heading 1}
\DefineParaStyle{Maple Warning}
\DefineParaStyle{Maple Heading 4}
\DefineParaStyle{Maple Heading 2}
\DefineParaStyle{Maple Heading 3}
\DefineParaStyle{Maple Dash Item}
\DefineParaStyle{Maple Error}
\DefineParaStyle{Maple Title}
\DefineParaStyle{Maple Ordered List 1}
\DefineParaStyle{Maple Text Output}
\DefineParaStyle{Maple Ordered List 2}
\DefineParaStyle{Maple Ordered List 3}
\DefineParaStyle{Maple Normal}
\DefineParaStyle{Maple Ordered List 4}
\DefineParaStyle{Maple Ordered List 5}
\DefineCharStyle{Maple 2D Output}
\DefineCharStyle{Maple 2D Input}
\DefineCharStyle{Maple Maple Input}
\DefineCharStyle{Maple 2D Math}
\DefineCharStyle{Maple Hyperlink}
\begin{center}
\begin{Maple Normal}
\underline{\textbf{Integration 1}}
\end{Maple Normal}
\end{center}
\begin{Maple Normal}
\_\_\_\_\_\_\_\_\_\_\_\_\_\_\_\_\_\_\_\_\_\_\_\_\_\_\_\_\_\_\_\_\_\_\_\_\_\_\_\_\_\_\_\_\_\_\_\_\_\_\_\_\_\_\_\_\_\_\_\_\_\_\_\_\_\_\_\_\_\_\_\_\_\_\_\_\_\_\_\_\_
\end{Maple Normal}
\section{\textbf{The }\textcolor[RGB]{255,0,0}{\textbf{ratint}}\textbf{ function}}
\begin{lstlisting}
> myarctan := 
 proc(y,x) 
  `if`(nargs = 1, arctan(y),evalf(Pi/2)-arctan(expand(x/y)));
 end:
myarctanh := (y) -> myln(1+y)/2 - myln(1-y)/2:
myln := 
 proc(u)
  local c,d,v;
  d := degree(u,x);
  c := coeff(u,x,d);
  v := expand(x^d + (u-c*x^d)/c);
  if d = 1 then v := abs(v); fi;
  ln(v);
 end:ratint := 
 proc(f,x)
  local g;
  g := int(f,x);
  if has(g,\{RootOf,int\}) then 
    g := int(convert(evalf(f),parfrac,x),x);
  fi;
  g := eval(
        subs(ln=myln,arctan = myarctan,arctanh=myarctanh,
             signum=1,abs=(x->x),
             g));
  g := expand(g);
  if (type(g,`+`)) then 
   g := remove((a)->type(a,constant),g);
  fi;
  g;
 end:
\end{lstlisting}
\section{\textbf{Various rational functions}}
\begin{lstlisting}
> g[1] := (x) -> x^2 + x + 1 + 1/x + 1/x^2;g[2] := (x) -> 1/(x-1) + 1/(x^2-1) + 1/(x^3-1);g[3] := (x) -> 1/(x-1)^2 + 1/(x-2)^3 + 1/(x-3)^4;g[4] := (x) -> 1/(x+1) + 1/(x^2+1) + 1/(x^3+1);g[5] := (x) -> x/(x+1)*(x+2)/(x+3)*(x+4);g[6] := (x) -> x/(x+1)^2*(x+2)/(x+3)^2*(x+4);g[7] := (x) -> (x^3+4*x^2+x+2)/(x-1)^2/(x+1)^3;g[8] := (x) -> 1/(x^2+8*x+1);g[9] := (x) -> (x^2+1)^(-5);g[10] := (x) -> 1/(x^8+x^2+1);
\end{lstlisting}
\begin{lstlisting}
> 
\end{lstlisting}
\begin{Maple Normal}
\_\_\_\_\_\_\_\_\_\_\_\_\_\_\_\_\_\_\_\_\_\_\_\_\_\_\_\_\_\_\_\_\_\_\_\_\_\_\_\_\_\_\_\_\_\_\_\_\_\_\_\_\_\_\_\_\_\_\_\_\_\_\_\_\_\_\_\_\_\_\_\_\_\_\_\_\_\_
\end{Maple Normal}
\begin{Maple Heading 2}
\textbf{Exercise 1}
\end{Maple Heading 2}
\begin{Maple Normal}
When 
$x$is large, 
$y$is approximately 
$\frac{x^{4}}{4}$.  Indeed, the formula for 
$y$ is as follows:
\end{Maple Normal}
\begin{lstlisting}
> y := ratint(((x^3+1)/(x^2+1))^3,x);
\end{lstlisting}
\begin{Maple Normal}
The first term is 
$\frac{x^{4}}{4}$, and the other terms are much smaller, at least when 
$x$ is large.  Indeed, when 
$x$\

is large, all of the terms with 
$x^{2}+1$ on the bottom will be very small.  Also, you should recall from\

the theory of trigonometric functions that 
${| \arctan (x)|}$ is always at most 
$\frac{\mathrm{Pi}}{2}$, so this will also be negligible\

in comparison to 
$\frac{x^{4}}{4}$ for large 
$x$.  The terms 
$3 x$ and 
$\frac{3 x^{2}}{2}$ will get large as 
$x$ does, but much more \

slowly than 
$\frac{x^{4}}{4}$ does.  For example, when 
$x =1000$, the term 
$\frac{x^{4}}{4}$ is of the order of a trillion, whereas \


$\frac{3 x^{2}}{2}$ is of the order of a million, and 
$3 x$ is only 
$3000$.  The behaviour of 
$\ln (x^{2}+1)$ is a little less obvious,\

but in fact it grows much more slowly even than 
$x$, as you could check by plotting the graph.\

\

We can ask Maple to separate the terms as follows:
\end{Maple Normal}
\begin{lstlisting}
> terms := [op(y)];
\end{lstlisting}
\begin{Maple Normal}
We can then check the relative size when 
$x =1000$:
\end{Maple Normal}
\begin{lstlisting}
> evalf(subs(x=1000,terms));
\end{lstlisting}
\begin{Maple Normal}

\end{Maple Normal}
\begin{Maple Normal}
You can see that the first term is much bigger than the rest.  We can also see this by plotting.\

In the picture below, the red curve is 
$\frac{x^{4}}{4}$, and the other terms are plotted in various different colours.
\end{Maple Normal}
\begin{Maple Normal}
By 
$x =5$, the red curve is already much bigger than the others.  If you change the range to \textcolor[RGB]{255,0,0}{\textbf{1..100}}\

(say) then only the red curve is visible because all other terms are squashed down onto the 
$x$ axis. 
\end{Maple Normal}
\begin{lstlisting}
> plot(terms,x=1..5);
\end{lstlisting}
\begin{Maple Normal}
We can also just plot 
$y$ and 
$\frac{x^{4}}{4}$ together.  If we stop at 
$x =9$ then you can still see the difference between\

the two curves, but if we went much further than that then they would be indistinguishable.
\end{Maple Normal}
\begin{lstlisting}
> plot([y,x^4/4],x=-9..9);
\end{lstlisting}
\begin{lstlisting}
> unassign('y','terms');
\end{lstlisting}
\begin{Maple Normal}
\_\_\_\_\_\_\_\_\_\_\_\_\_\_\_\_\_\_\_\_\_\_\_\_\_\_\_\_\_\_\_\_\_\_\_\_\_\_\_\_\_\_\_\_\_\_\_\_\_\_\_\_\_\_\_\_\_\_\_\_\_\_\_\_\_\_\_\_\_\_\_\_\_\_\_\_\_\_\_\_\_
\end{Maple Normal}
\begin{Maple Heading 2}
\textbf{Exercise 2}
\end{Maple Heading 2}
\begin{Maple Normal}
Here is the formula for 
$y$:
\end{Maple Normal}
\begin{lstlisting}
> y := ratint((a*x^2+b)/(c*x^2+d),x);
\end{lstlisting}
\begin{Maple Normal}
This is 
$y =\frac{a x}{c}$ (which gives a line of slope 
$\frac{a}{c}$) plus some other terms, which will be much smaller when 
$x$ is large.  One way to see this is to look at  
$\frac{\mathit{dy}}{\mathit{dx}}$ when 
$x$ is large, or in other words, the limit of 
$\frac{\mathit{dy}}{\mathit{dx}}$ as 
$x$ tends to infinity.  Of course 
$\frac{\mathit{dy}}{\mathit{dx}}$ is just the function we first thought of, namely 
$\frac{a x^{2}+b}{c x^{2}+d}$.  When 
$x$ is very large, 
$b$ and 
$d$ will be negligible compared with 
$a x^{2}$ and 
$c x^{2}$, so the function is approximately 
$\frac{a x^{2}}{c x^{2}}=\frac{a}{c}$.
\end{Maple Normal}
\begin{lstlisting}
> limit(diff(y,x),x=infinity);
\end{lstlisting}
\begin{Maple Normal}

\end{Maple Normal}
\begin{Maple Normal}
We can see this graphically if we choose some numbers for 
$a ,b ,c$ and 
$d$.
\end{Maple Normal}
\begin{lstlisting}
> plot(subs(a=7,b=-3,c=2,d=8,[y,a*x/c]),x=10..100);
\end{lstlisting}
\begin{lstlisting}
> unassign('y');
\end{lstlisting}
\begin{Maple Normal}
\_\_\_\_\_\_\_\_\_\_\_\_\_\_\_\_\_\_\_\_\_\_\_\_\_\_\_\_\_\_\_\_\_\_\_\_\_\_\_\_\_\_\_\_\_\_\_\_\_\_\_\_\_\_\_\_\_\_\_\_\_\_\_\_\_\_\_\_\_\_\_\_\_\_\_\_\_\_\_\_\_
\end{Maple Normal}
\begin{Maple Heading 2}
\textbf{Exercise 3}
\end{Maple Heading 2}
\begin{Maple Normal}
\textbf{(a)} It is true that the integral of a rational function can contain terms like 
$a \ln (x^{2}+u x +v)$.  For example:
\end{Maple Normal}
\begin{lstlisting}
> y := (8*x+8)/(x^2+2*x+3);
\end{lstlisting}
\begin{lstlisting}
> ratint(y,x);
\end{lstlisting}
\begin{Maple Normal}
or more generally:
\end{Maple Normal}
\begin{lstlisting}
> y := a*(2*x+u)/(x^2+u*x+v);
\end{lstlisting}
\begin{lstlisting}
> ratint(y,x);
\end{lstlisting}
\begin{Maple Normal}

\end{Maple Normal}
\begin{Maple Normal}

\end{Maple Normal}
\begin{Maple Normal}
\textbf{(b)} It is \textit{not} true that terms like 
$x \ln (x +u)$ can occur in the integral of a rational function 
$g (x)$.  
\end{Maple Normal}
\begin{Maple Normal}
If there were a term 
$x \ln (x +u)$ in 
$\int g (x)d x$, then there would have to be a term 
\end{Maple Normal}
\begin{Maple Normal}

$\frac{\partial}{\partial x} (x \ln (x +u))$ in 
$\frac{d}{d x} (\int g (x)d x)$.  In other words, there would be a term 
\end{Maple Normal}
\begin{Maple Normal}

$\ln (x +u)+\frac{x}{x +u}$ in 
$g (x)$, which is not allowed, as 
$g (x)$ is supposed to be a rational function.
\end{Maple Normal}
\begin{Maple Normal}

\end{Maple Normal}
\begin{Maple Normal}

\end{Maple Normal}
\begin{Maple Normal}
\textbf{(c) }For essentially the same reason, there cannot be any terms like 
$a \ln (x +u)^{2}$ in the integral 
\end{Maple Normal}
\begin{Maple Normal}
of a rational function, because 
$\frac{\partial}{\partial x} (a \ln (x +u)^{2})=\frac{2 a \ln (x +u)}{x -u}$, which is not allowed as
\end{Maple Normal}
\begin{Maple Normal}
a term in a rational function.
\end{Maple Normal}
\begin{Maple Normal}
\_\_\_\_\_\_\_\_\_\_\_\_\_\_\_\_\_\_\_\_\_\_\_\_\_\_\_\_\_\_\_\_\_\_\_\_\_\_\_\_\_\_\_\_\_\_\_\_\_\_\_\_\_\_\_\_\_\_\_\_\_\_\_\_\_\_\_\_\_\_\_\_\_\_\_\_\_\_\_\_\_
\end{Maple Normal}
\begin{Maple Heading 2}
\textbf{Exercise 4}
\end{Maple Heading 2}
\begin{Maple Normal}
If 
$b^{2}-4 c <0$ then the function 
$x^{2}+b x +c$ has no real roots, and the integral 
$\int \frac{1}{x^{2}+b x +c}d x$ involves
\end{Maple Normal}
\begin{Maple Normal}
the 
$\arctan ()$ function.   Here is an example:
\end{Maple Normal}
\begin{lstlisting}
> b := 2; c := 3; b^2-4*c;
\end{lstlisting}
\begin{lstlisting}
> plot(x^2+b*x+c,x=-2..1,0..6);

\end{lstlisting}
\begin{lstlisting}
> ratint(1/(x^2+b*x+c),x);
\end{lstlisting}
\begin{Maple Normal}
If 
$b^{2}-4 c =0$ then the function 
$x^{2}+b x +c$ has only one real root (where the graph 
\end{Maple Normal}
\begin{Maple Normal}
touches the 
$x$ axis but does not cross it) and the integral involves neither 
$\arctan ()$ nor 
$\ln ()$; in fact, it just
\end{Maple Normal}
\begin{Maple Normal}
has the form 
$-\frac{1}{x -u}$.  Here is an example:
\end{Maple Normal}
\begin{lstlisting}
> b := 4; c := 4; b^2-4*c;
\end{lstlisting}
\begin{lstlisting}
> plot(x^2+b*x+c,x=-4..1,0..6);

\end{lstlisting}
\begin{lstlisting}
> ratint(1/(x^2+b*x+c),x);
\end{lstlisting}
\begin{Maple Normal}
If 
$0<b^{2}-4 c$ then the function 
$x^{2}+b x +c$ has two real roots,  and the integral is a
\end{Maple Normal}
\begin{Maple Normal}
sum of two terms involving 
$\ln ()$.  Here is an example:
\end{Maple Normal}
\begin{lstlisting}
> b := -5; c := 4; b^2-4*c;
\end{lstlisting}
\begin{lstlisting}
> plot(x^2+b*x+c,x=-1..5);

\end{lstlisting}
\begin{lstlisting}
> ratint(1/(x^2+b*x+c),x);
\end{lstlisting}
\begin{lstlisting}
> unassign('b','c');
\end{lstlisting}
\begin{Maple Normal}
\_\_\_\_\_\_\_\_\_\_\_\_\_\_\_\_\_\_\_\_\_\_\_\_\_\_\_\_\_\_\_\_\_\_\_\_\_\_\_\_\_\_\_\_\_\_\_\_\_\_\_\_\_\_\_\_\_\_\_\_\_\_\_\_\_\_\_\_\_\_\_\_\_\_\_\_\_\_\_\_\_
\end{Maple Normal}
\begin{Maple Heading 2}
\textbf{Exercise 5}
\end{Maple Heading 2}
\begin{Maple Normal}
Terms like 
$\ln ({| x -u |})$ in 
$\textcolor{gray}{\int}g (x)\textcolor{gray}{d}x$ are associated with places where 
$g (x)$ blows up to infinity (which
\end{Maple Normal}
\begin{Maple Normal}
are known as the \textit{poles }of 
$g (x)$).  More precisely, if there is a term 
$\ln ({| x -u |})$ in the integral, then
\end{Maple Normal}
\begin{Maple Normal}

$g (x)$ \textit{must }blow up at 
$x =u$.  Conversely, if 
$g (x)$ blows up at 
$x =u$ then there will \textit{usually} be a term
\end{Maple Normal}
\begin{Maple Normal}

$\ln ({| x -u |})$ in the integral, but not always; there may instead be terms of the form 
$(x -u)^{-n}$.  The
\end{Maple Normal}
\begin{Maple Normal}
functions 
$g_{3}(x)$ and 
$g_{7}(x)$ are examples of this.
\end{Maple Normal}
\begin{Maple Normal}
  
\end{Maple Normal}
\begin{Maple Normal}

\end{Maple Normal}
\begin{Maple Normal}
Here we consider 
$g_{2}(x)$; there are poles at 
$x =-1$ and 
$x =1$, and corresponding terms 
$\ln ({| x +1|})$
\end{Maple Normal}
\begin{Maple Normal}
and 
$\ln ({| x -1|})$ in the integral.
\end{Maple Normal}
\begin{lstlisting}
> y := g[2](x); z := ratint(y,x);
\end{lstlisting}
\begin{lstlisting}
> plot([y,z],x=-5..5,-10..10);
\end{lstlisting}
\begin{Maple Normal}

\end{Maple Normal}
\begin{Maple Normal}

\end{Maple Normal}
\begin{Maple Normal}
Here we consider 
$g_{3}(x)$.   There are no 
$\ln ()$ terms in the integral, but nonetheless there are poles
\end{Maple Normal}
\begin{Maple Normal}
at 
$x =1$, 
$x =2$ and 
$x =3$.
\end{Maple Normal}
\begin{lstlisting}
> y := g[3](x); z := ratint(y,x);
\end{lstlisting}
\begin{lstlisting}
> plot([y,z],x=0..5,-100..100,numpoints=400);
\end{lstlisting}
\begin{Maple Normal}

\end{Maple Normal}
\begin{Maple Normal}
Here we consider 
$g_{4}(x)$.  There is a pole at 
$x =-1$, and a corresponding term 
$\ln ({| x +1|})$ in the
\end{Maple Normal}
\begin{Maple Normal}
integral.
\end{Maple Normal}
\begin{lstlisting}
> y := g[4](x); z := ratint(y,x);
\end{lstlisting}
\begin{lstlisting}
> plot([y,z],x=-5..5,-10..10);
\end{lstlisting}
\begin{Maple Normal}
\_\_\_\_\_\_\_\_\_\_\_\_\_\_\_\_\_\_\_\_\_\_\_\_\_\_\_\_\_\_\_\_\_\_\_\_\_\_\_\_\_\_\_\_\_\_\_\_\_\_\_\_\_\_\_\_\_\_\_\_\_\_\_\_\_\_\_\_\_\_\_\_\_\_\_\_\_\_\_\_\_
\end{Maple Normal}
\begin{Maple Heading 2}
\textbf{Exercise 6}
\end{Maple Heading 2}
\begin{Maple Normal}
We only get terms like 
$\frac{1}{x -u}$ or 
$\frac{1}{(x -u)^{n}}$ in 
$\textcolor{gray}{\int}g (x)\textcolor{gray}{d}x$ if the denominator of 
$g (x)$ has repeated roots.
\end{Maple Normal}
\begin{Maple Normal}

\end{Maple Normal}
\begin{Maple Normal}
A repeated root at 
$x =u$ gives a factor 
$(x -u)^{n}$ in the denominator, with 
$1<n$.  This gives a term 
$\frac{1}{(x -u)^{n -1}}$ in the integral (and possibly also some terms like 
$\frac{1}{(x -u)^{k}}$ for smaller values of 
$k$).
\end{Maple Normal}
\begin{Maple Normal}

\end{Maple Normal}
\begin{Maple Normal}
Here we consider 
$g_{3}(x)$:
\end{Maple Normal}
\begin{lstlisting}
> y := g[3](x);
\end{lstlisting}
\begin{Maple Normal}
The denominator has repeated roots at 
$x =1$, 
$x =2$ and 
$x =3$:
\end{Maple Normal}
\begin{lstlisting}
> d := denom(factor(y));
\end{lstlisting}
\begin{Maple Normal}
These show up in the graph of 
$d$ as points where the graph touches the 
$x$ axis without crossing it,\

or where it crosses but the tangent line at the crossing point is horizontal: 
\end{Maple Normal}
\begin{lstlisting}
> plot(d,x=0..4,-0.3...0.3);
\end{lstlisting}
\begin{Maple Normal}
The integral has terms of the form 
$\frac{1}{(x -u)^{n}}$ for 
$u =1,2$ and 
$3$.
\end{Maple Normal}
\begin{lstlisting}
> ratint(y,x);
\end{lstlisting}
\begin{Maple Normal}

\end{Maple Normal}
\begin{Maple Normal}

\end{Maple Normal}
\begin{Maple Normal}

\end{Maple Normal}
\begin{Maple Normal}
Here we consider 
$g_{4}(x)$:
\end{Maple Normal}
\begin{lstlisting}
> y := g[4](x);
\end{lstlisting}
\begin{Maple Normal}
The denominator has a root at 
$x =-1$, which is not repeated:
\end{Maple Normal}
\begin{lstlisting}
> d := denom(factor(y));
\end{lstlisting}
\begin{Maple Normal}
This is visible in the graph of 
$d$.  Note that where the graph meets the 
$x$ axis, the tangent line is not
\end{Maple Normal}
\begin{Maple Normal}
horizontal, corresponding to the fact that the root is not repeated.
\end{Maple Normal}
\begin{lstlisting}
> plot(d,x=-2..2,-10..10);
\end{lstlisting}
\begin{Maple Normal}
The integral has no terms of the form 
$\frac{1}{(x -u)^{n}}$.
\end{Maple Normal}
\begin{lstlisting}
> ratint(y,x);
\end{lstlisting}
\begin{lstlisting}
> 
\end{lstlisting}
\begin{Maple Normal}

\end{Maple Normal}
\begin{Maple Normal}

\end{Maple Normal}
\begin{Maple Normal}
Here we consider 
$g_{6}(x)$:
\end{Maple Normal}
\begin{lstlisting}
> y := g[6](x);
\end{lstlisting}
\begin{Maple Normal}
The denominator has repeated roots at 
$x =-3$ and 
$x =-1$:
\end{Maple Normal}
\begin{lstlisting}
> d := denom(factor(y));
\end{lstlisting}
\begin{Maple Normal}
These visible in the graph of 
$d$, as points where the curve meets the 
$x$ axis without crossing it.
\end{Maple Normal}
\begin{Maple Normal}
Note that the tangent line is horizontal at these points.
\end{Maple Normal}
\begin{lstlisting}
> plot(d,x=-4..2,-10..10);
\end{lstlisting}
\begin{Maple Normal}
The integral has terms of the form 
$\frac{1}{x -u}$ for 
$x =-3$ and 
$x =-1$.
\end{Maple Normal}
\begin{lstlisting}
> ratint(y,x);
\end{lstlisting}
\begin{Maple Normal}

\end{Maple Normal}
\begin{Maple Normal}

\end{Maple Normal}
\begin{Maple Normal}
We now integrate some random rational functions, and observe that we do not get any terms like
\end{Maple Normal}
\begin{Maple Normal}

$\frac{1}{(x -u)^{n}}$.  The reason is simply that the denominator is a random polynomial, which will almost
\end{Maple Normal}
\begin{Maple Normal}
never have repeated roots.  The roots of a random polynomial are randomly distributed, so it 
\end{Maple Normal}
\begin{Maple Normal}
would be a big coincidence for two roots to be in the same place, and this does not often happen.
\end{Maple Normal}
\begin{lstlisting}
> r := randpoly(x)/randpoly(x);
\end{lstlisting}
\begin{lstlisting}
> ratint(r,x);
\end{lstlisting}
\begin{lstlisting}
> 
\end{lstlisting}
\begin{lstlisting}
> r := randpoly(x)/randpoly(x);
\end{lstlisting}
\begin{lstlisting}
> ratint(r,x);
\end{lstlisting}
\begin{lstlisting}
> 
\end{lstlisting}
\begin{lstlisting}
> r := randpoly(x)/randpoly(x);
\end{lstlisting}
\begin{lstlisting}
> ratint(r,x);
\end{lstlisting}
\begin{lstlisting}
> 
\end{lstlisting}
\end{document}
